\section{Wie Materie zur Quelle der Geometrie werden könnte}
\label{sec:geometry}

Angenommen, das Modell könnte den ausgewählten Materiezuständen Massen zuordnen. Wir wissen dann noch nicht, wie dieselben Zustände auf die Geometrie wirken. Dazu braucht es eine Quelle und eine Gleichung für ihre Wirkung. Der erste Test ist, ob die ausgelesenen Größen eine geschlossene Dynamik haben. Erst dann stellt sich die Frage, welche Feldgleichung diese Dynamik beschreiben kann.

\subsection{Geschlossene Dynamik ausgelesener Variablen}

Die Zustandszerlegung aus Abschnitt \ref{sec:carrier} lässt sich auf eine zeitabhängige Auslese erweitern. Sei \(u\) ein Zustand eines endlichdimensionalen reellen Vergleichssystems mit \(\dot u=G(t)u\). Die Auslese sei \(r=T(t)u\) mit differenzierbarem surjektivem \(T(t)\). Zu einer Rechtsinversen \(R(t)\) von \(T(t)\) lässt sich jeder Zustand als \(u=Rr+z\) mit \(Tz=0\) schreiben. Dann gilt
\begin{align}
 \dot r&=(\dot T+TG)Rr+(\dot T+TG)z,\\
 \dot r\ \text{ist durch }(t,r)\text{ bestimmt}
      &\quad\Longleftrightarrow\quad (\dot T+TG)\ker T=0 .
 \label{eq:autonomy}
\end{align}
Die Notwendigkeit folgt, indem zu einem Zustand ein beliebiges Element des Kerns addiert wird. Die Hinreichendheit folgt aus dem Verschwinden des zweiten Terms. Der erste Term ist dann unabhängig von der Wahl der Rechtsinversen. Für ein konstantes \(T\) verlangt \eqref{eq:autonomy}, dass \(G\) den Kern erhält. Für eine zeitabhängige Auslese gehört ihre eigene Änderung zur Bilanz \cite{physics}.

Das Resultat betrifft einen linearen endlichen Vergleich. In einer Feldtheorie kommen unter anderem Definitionsbereiche von Differentialoperatoren, Randbedingungen, Eichfreiheit und die Fortpflanzung von Zwangsbedingungen hinzu. \eqref{eq:autonomy} macht die lokale Frage nach den ausgelassenen Zuständen sichtbar. Es entscheidet diese weiteren Fragen nicht.

\paragraph{Was ein verschwundener Anteil noch bewirken kann.}
Für die feste HC-Projektion aus \eqref{eq:projections} ist \(r=Cu\). Dann führt \eqref{eq:autonomy} auf \(CGP_-=0\). Als Beispiel vertausche \(J\) im Zustandsraum \(\mathbb R^2\) die beiden Koordinaten und sei \(G=\operatorname{diag}(1,0)\). Der Zustand \(u=(1,-1)^{\mathsf T}\) hat \(Cu=0\) und dennoch \(CGu=(1,1)^{\mathsf T}/2\). Aus dem ausgelesenen Zustand allein kann man also in diesem Vergleich seine Ableitung nicht bestimmen \cite{physics}.

\subsection{Coframe und Quelle in einem Cartan-Vergleich}

Eine mögliche geometrische Auslese ist ein reeller, punktweise invertierbarer Coframe \(E^a{}_\mu\). Dann lässt sich
\begin{equation}
 g_{\mu\nu}=\eta_{ab}E^a{}_\mu E^b{}_\nu,\qquad
 \eta=\operatorname{diag}(-1,1,1,1)
 \label{eq:metric}
\end{equation}
als Lorentzmetrik definieren. Eine bilineare Auslese für den Coframe muss daher vollen Rang haben.

Der explizite Einstein--Cartan-Vergleich des Companion setzt Coframe, Lorentzverbindung und ein Spinortripel als unabhängige Felder an. Ein Multiplikator koppelt den Coframe an eine bilineare Spinorauslese \(\mathcal C^a{}_\mu(\Psi,D\Psi)\). Die vorgeschriebene Wirkung hat die Form
\begin{equation}
 S_{\mathrm{vgl}}=\int\! d^4x\,
 \left[\left|\det E\right|\left(\tfrac c2 R(E,\omega)-u_0\right)
 +\mathcal L_{\mathrm D}
 +\lambda_a{}^\mu\bigl(E^a{}_\mu-\mathcal C^a{}_\mu\bigr)\right].
 \label{eq:cartan}
\end{equation}
Hier ist \(\mathcal L_{\mathrm D}\) die dort definierte Spinor-Lagrangedichte, \(\omega\) die unabhängige Lorentzverbindung, \(c\) ein Vergleichskoeffizient und \(u_0\) ein konstanter Term. Die Variation nach \(\lambda_a{}^\mu\) liefert \(E^a{}_\mu=\mathcal C^a{}_\mu\). Die Variationen nach Spinoren, Verbindung und Coframe stellen weitere Gleichungen. Insbesondere muss eine algebraisch bestimmte Verbindung mit ihren tatsächlich daraus folgenden Ableitungen eingesetzt werden, wenn man eine reduzierte Feldformulierung prüft \cite{companion,repo}.

In der konkreten neutralen Spinorrechnung wirken die beiden Vergleichsmatrizen \(Q\) und \(M_{\mathrm D}\) auf die drei Kanäle. Mit \(v=|\det E|\), \(e=E^{-1}\), \(K^a=\beta\Gamma^a\) und \(D_\mu\Psi=\partial_\mu\Psi-\tfrac14\omega_{\mu,ab}\Gamma^a\Gamma^b\Psi\) werden die Auslese und die Dichte durch
\begin{align}
 \mathcal C^a{}_\mu&=\operatorname{Im}(\Psi^\dagger QK^aD_\mu\Psi),\\
 \mathcal L_{\mathrm D}
 &=-v\operatorname{Im}(\Psi^\dagger e_a{}^\mu K^aD_\mu\Psi)
   -v\operatorname{Re}(\Psi^\dagger(M_{\mathrm D}\otimes\beta)\Psi)
\end{align}
festgelegt. Die Matrizen \(\Gamma^a,\beta\) bezeichnen die Dirac-Matrizen in der Vergleichskonvention. Kanal- und Spinindizes werden getrennt kontrahiert. Diese Festlegung macht sichtbar, welche Spinorgröße für die geometrische Auslese eingesetzt wird und an welcher Stelle ein Massenoperator in die Spinordynamik eingeht.

Die Vergleichswirkung setzt ihre Krümmung, Materiedichte und Kopplungskoeffizienten als Modellbestandteile an. Sie zeigt, wo eine Mate\-rie\-größe und der Coframe in derselben Variationsrechnung zusammentreffen. Für den nativen Übergang aus Abschnitt \ref{sec:mass} wäre zusätzlich zu bestimmen, welche dieser Eingaben aus dem gemeinsamen Träger folgen und wie die physikalische Quelle an die dortige Massenauslese gekoppelt ist.

\subsection{Von Daten an einem Punkt zu einer Feldlösung}

Die Cartan-Rechnungen im Companion prüfen endliche Ableitungsdaten an einem Punkt. Nach Eliminierung höherer freier Ableitungen ergeben sich lineare Bedingungen an vorgeschriebene Daten und den Quellenzuwachs. Ist die reduzierte Quellenmatrix \(L\) von vollem Spaltenrang, besitzt \(Lq=b\) für kompatible Daten genau eine Lösung \(q=L^+b\), wobei \(L^+\) die Moore--Penrose-Pseudoinverse ist. Die Kompatibilitätsbedingung lautet \((I-LL^+)b=0\). Diese Rechnung kann angeben, welche Daten an diesem Punkt eine gemeinsame Fortsetzung der geprüften endlichen Ordnung zulassen.

Der Vergleich enthält bereits eine bedingte Aussage zur Fortpflanzung der gravitativen Zwangsbedingungen. Sind die Verbindung kompatibel eliminiert und sowohl die Konstitutivgleichung \(E^a{}_\mu=\mathcal C^a{}_\mu\) aus der Variation nach \(\lambda_a{}^\mu\) als auch die Sondengleichung, die Euler-Gleichung der Spinoren, in einem offenen Gebiet erfüllt, geben die Ward-Identitäten dort einen symmetrischen und divergenzfreien Coframe-Resttensor. Mit den räumlich projizierten spurmodifizierten Coframe-Gleichungen erhält man unter den angegebenen Regularitäts- und Randbedingungen ein symmetrisch hyperbolisches System für die vier gravitativen Zwangsgrößen. Bei verschwindenden Anfangsdaten und ohne einlaufenden Zwangsfluss bleiben diese Größen null. Die Anwendung setzt die beiden zuerst genannten Feldgleichungen auf dem Gebiet voraus. Ein Verschwinden ihrer Reste an einem einzelnen Punkt reicht dafür nicht aus \cite{companion,repo}.

Eine Lösung in einem offenen Gebiet erfordert zusätzlich die gleichzeitige Erfüllung der Konstitutiv- und Sondengleichungen, kompatible Flächendaten und eine Kontrolle der Fortpflanzung der Zwangsbedingungen. Das bleibt eine eigene Aufgabe. Entsprechend unterscheidet der Companion die punktweise Prüfung des Cartan-Vergleichs von der allgemeinen Schließung des Feldsystems.

\subsection{Vom Materieantrieb zur geometrischen Antwort}
\label{sec:op5-coupling}

Im Cartan-Vergleich tritt ein vorgegebener Massenterm in der Spinor-Lagrangedichte auf. Für die Vereinheitlichungsfrage wollen wir weitergehen. Wie kann eine Änderung der Materie eine geometrische Antwort erzeugen, und woran lässt sie sich von der gewöhnlichen Gravitation des Aufbaus unterscheiden? Dazu verfolgen wir ein und dieselbe Modulation von der Materiequelle bis zur Auslese. Die Rechnung spielt im Materieframe \(M\) und benutzt Einheiten mit Lichtgeschwindigkeit eins. Als Eingabe dienen die Massenparameter des Dirac--Yukawa-Vergleichs. Die Identifikation dieser Parameter mit der nativen Massenauslese aus Abschnitt \ref{sec:mass} ist weiterhin offen.

Die beiden Wege in der folgenden Skizze setzen bei derselben Materialvariation an. Über den Energie-Impuls-Tensor ist ihr gewöhnlicher Beitrag bestimmt, sobald die mechanische Quelle und die Randdaten festliegen. Für den gerichteten Beitrag kommt zusätzlich die Zuordnung der radialen Higgsquelle zur Quelle des Vergleichsfelds hinzu. Genau diese Zuordnung wird unten als offene Annahme sichtbar.

\begin{figure}[!htbp]
\centering
\begin{tikzpicture}[>=Latex,font=\small,
  box/.style={draw=qs axis,rounded corners=2pt,align=center,text width=51mm,
             minimum height=10mm,inner sep=4pt,fill=white}]
  \node[box,text width=65mm] (density) at (0,0) {Präparierte Materialverschiebung\\\(\delta n_i\) und \(\delta T_{\mu\nu}\)};
  \node[box,fill=qs blue!7] (gravity) at (-3.55,-1.61)
     {Gewöhnliche Gravitation\\\(D_0\bar\gamma_{\mu\nu}=16\pi G\,\delta T_{\mu\nu}\)};
  \node[box,fill=qs rust!9] (higgs) at (3.55,-1.61)
     {Radiale Higgsquelle\\\(\delta\mathcal J_h\) aus W2};
  \node[box,fill=qs blue!7] (exterior) at (-3.55,-4.13)
     {Kompakter linearer\\mechanischer Beitrag};
  \node[box,fill=qs rust!9] (profile) at (3.55,-4.13)
     {Gerichtetes OP5-Profil\\\(\delta\mathcal J_{\rm OP5}=\mathcal P_{\rm dyn}\delta H\)};
  \node[box,fill=qs rust!9] (axis) at (3.55,-5.66)
     {Gezeitenkontrast auf dem Plateau\\\(\Delta\mathcal I\)};
  \draw[->,thick,qs blue] (density.south west) -- (gravity.north);
  \draw[->,thick,qs rust] (density.south east) -- (higgs.north);
  \draw[->,thick,qs blue] (gravity) -- (exterior);
  \draw[->,thick,dashed,qs rust] (higgs) -- (profile);
  \node[anchor=east,align=right,inner sep=4pt] at (3.52,-2.9)
      {angesetzte Quellenidentifikation\\[-1pt]\(\delta\mathcal J_{\rm OP5}=\delta\mathcal J_h\)};
  \draw[->,thick,qs rust] (profile) -- (axis);
\end{tikzpicture}
\caption{Die beiden Antwortwege derselben linearen Materialpräparation. Der gewöhnliche Beitrag folgt aus dem mechanischen Tensor. Der gerichtete Vergleich setzt die Kopplung der radialen W2-Quelle an das OP5-Profil voraus. Die gestrichelte Verbindung markiert diese Annahme.}
\label{fig:op5-modell}
\end{figure}


\subsubsection{Welche Quelle die Massenterme liefern}

W2 bezeichnet die gemeinsame Vergleichswirkung des Companion: Gravitation nach Einstein und Cartan mit den Spinoren einer Generation und ihren Eichfeldern (W1), ergänzt um ein Higgs-Dublett mit Yukawa-Kopplungen (Runden v4.9.0-r81 und v4.9.0-r82 \cite{companion}). Sei \(h_{\rm W}=v+\eta_h\) der Higgsradius von W2 bei festen Higgswinkeln und \(v>0\) sein Referenzwert. In der gewählten Normierung gilt \(m_i(h_{\rm W})=y_i h_{\rm W}\). Für nichtrelativistische Besetzungen mit Teilchenzahldichten \(n_i\) behalten wir zunächst die Ruheenergie. Dabei setzen wir die Ruheenergie des Materials mit seiner Yukawa-Masse gleich, nehmen also an, dass sie ganz vom Higgsradius getragen wird. Das ist eine Modellannahme dieser Fallstudie; Runde v4.10.0-r35 \cite{layers} führt den Higgsanteil der Ruhemasse als eigene Größe. Die Materiedichte und ihre radiale Quelle sind dann
\begin{align}
 \epsilon_{\rm mat}(h_{\rm W},n_i)&=h_{\rm W}\sum_i y_i n_i,
 &
 \mathcal J_h&=-\partial_{h_{\rm W}}\epsilon_{\rm mat}
              =-\sum_i y_i n_i .
 \label{eq:op5-yukawa}
\end{align}
Die Dichte \(n_i\) zählt bei beidseitiger Besetzung Teilchen und Antiteilchen gemeinsam. Die dem Higgsradius zugeordnete Quelle hängt dabei von der Besetzungssumme ab. Eine kleine Verschiebung \(\mathbf u\) einer homogenen Referenz ändert diese durch \(\delta n_i=-n_{i0}\nabla\cdot\mathbf u\). Deshalb ist
\begin{equation}
 \delta\mathcal J_h=\left(\sum_i y_i n_{i0}\right)\nabla\cdot\mathbf u,
 \qquad
 \left.\delta\epsilon_{\rm mat}\right|_v=-v\,\delta\mathcal J_h .
 \label{eq:op5-density}
\end{equation}
Kompression und Verdünnung liefern Quellenstörungen mit entgegengesetztem Vorzeichen. Sie ändern zugleich die gewöhnliche Energiedichte. Bewegt sich auch der Higgsradius, kommt zu \(\left.\delta\epsilon_{\rm mat}\right|_v\) der Beitrag \(\eta_h\sum_i y_i n_{i0}\) hinzu. Der folgende mechanische Abgleich verfolgt den direkten Transport der Ruhemasse bei festem Referenzradius. Die Feldenergie und der Impulsaustausch mit einem veränderlichen Higgsfeld gehören zur vollständigen Bilanz \cite{op5-density}.

Die Quelle in \eqref{eq:op5-density} beschreibt zunächst einen Dichteantrieb. Das lässt sich auch am Druck prüfen. Bei fester Entropie pro Teilchen gilt für eine einzelne Art \(p_{\rm mat}=n\partial_n\epsilon_{\rm mat}-\epsilon_{\rm mat}\). Die reine Ruheenergie ist proportional zu \(n\), sodass ihr Beitrag zu \(p_{\rm mat}\) verschwindet. Elastische und kinetische Terme sind deshalb als eigene Beiträge zu führen. Der lokale Energieabgleich bei festem \(v\) ersetzt die Kraftbilanz eines bewegten Materials nicht. Mit der Signatur \((-+++)\) und der Konvention \(T_{\mu\nu}=-2\,\delta S_{\rm mat}/(\sqrt{-g}\,\delta g^{\mu\nu})\) ergibt die Diffeomorphismusidentität auf den Materiegleichungen
\begin{equation}
 \nabla_\mu T^\mu{}_\nu=\mathcal J_h\,\partial_\nu h_{\rm W}
 \label{eq:op5-matter-ward}
\end{equation}
für einen als extern behandelten Higgsradius und ohne weitere äußere Kräfte. Bei dynamischem \(h_{\rm W}\) ist diese rechte Seite Teil eines Austauschs, den die übrigen Felder in der Gesamtbilanz aufnehmen müssen. Im mechanischen Vergleich unten werden der Referenzradius festgehalten und die inneren elastischen Kräfte gemeinsam bilanziert. Damit lässt sich \eqref{eq:op5-matter-ward} von der Erhaltung des präparierten linearen Tensors abgrenzen \cite{op5-density}.

\subsubsection{Wie sich der gewöhnliche Außenbeitrag kontrollieren lässt}

Kann man diese Quelle schalten und ihren gewöhnlichen gravitativen Beitrag zugleich kontrollieren? Der Vergleich benutzt zwei überlagerte isotrope elastische Kontinua mit Ruhedichten \(\rho_1,\rho_2>0\). Der lineare Cauchy-Tensor der ersten Komponente ist \(\sigma(\mathbf u)=\mathsf C\,\varepsilon(\mathbf u)\). Die zweite Komponente hat die Steifigkeit \(c_\sigma\mathsf C\) mit \(c_\sigma>0\). Wir schreiben \(\mathcal L\mathbf u=\nabla\cdot\sigma(\mathbf u)\) und
\begin{equation}
 A_1=\rho_1\partial_t^2-\mathcal L,\qquad
 A_2=\rho_2\partial_t^2-c_\sigma\mathcal L .
 \label{eq:op5-elastic}
\end{equation}
Bei konstanten Koeffizienten kommutieren \(A_1\) und \(A_2\). Für ein glattes, räumlich kompaktes Vektorfeld \(\mathbf b\) wähle
\[
 \mathbf u=A_2\mathbf b,\quad
 \mathbf w=-A_1\mathbf b,\quad
 \mathbf d=\mathbf u-\mathbf w+k^{-1}A_1\mathbf u .
\]
Eine innere Paarfeder mit \(k>0\) und vorgeschriebener relativer Ruhelage \(\mathbf d\) liefert die entgegengesetzten Kräfte \(\mp k(\mathbf u-\mathbf w-\mathbf d)\). Somit gelten während der Schaltung beide Bewegungsgleichungen \(A_1\mathbf u=-k(\mathbf u-\mathbf w-\mathbf d)\) und \(A_2\mathbf w=+k(\mathbf u-\mathbf w-\mathbf d)\). Die gezielte Änderung von \(\mathbf d\) ist eine Eingabe des Vergleichs \cite{op5-motion}.

Die beiden Gleichungen für die Feder lassen sich direkt an der gewählten Verschiebung kontrollieren. Es gilt \(A_1\mathbf u=A_1A_2\mathbf b\) und \(A_2\mathbf w=-A_2A_1\mathbf b=-A_1\mathbf u\). Beim Einsetzen von \(\mathbf d\) wird die Federkraft auf die erste Komponente mit \(A_1\mathbf u\) identifiziert. Damit ist die Summe ihrer Kräfte im Referenzmodell zu jedem Zeitpunkt null. Auf dem statischen Plateau ist \(\mathbf w=-\mathbf u/c_\sigma\). Würde man dieses Verhältnis auch während der Schaltung vorgeben, verschwände die Summenspannung linear. Aus der Summe der beiden Bewegungsgleichungen folgte dann
\[
 (\rho_1-\rho_2/c_\sigma)\,\ddot{\mathbf u}=0.
\]
Mit \(\beta_\rho=\rho_2-c_\sigma\rho_1\) ist der Vorfaktor \(-\beta_\rho/c_\sigma\). Gerade bei \(\beta_\rho\ne0\) könnte dieser starre Ansatz also keinen Einschaltvorgang aus dem Ruhezustand tragen. Die zwei Operatoren in \eqref{eq:op5-elastic} lösen dieses Problem, indem die Summenspannung während der Schaltung die Änderung des Gesamtimpulses trägt \cite{op5-motion}.

Für die gewählte Verschiebung lauten die zugehörigen linearen mechanischen Beiträge
\begin{equation}
 \delta T^{00}=-\beta_\rho\nabla\cdot\mathcal L\mathbf b,\qquad
 \delta T^{0i}=\beta_\rho(\mathcal L\dot{\mathbf b})_i,\qquad
 \delta T^{ij}=-\beta_\rho\sigma_{ij}(\ddot{\mathbf b}) .
 \label{eq:op5-tensor}
\end{equation}
Zur Kontrolle kann man Dichte, Impuls und Spannung aus zwei verschiedenen Linearkombinationen der Verschiebungen lesen:
\begin{align}
 \rho_1\mathbf u+\rho_2\mathbf w
    &=\beta_\rho\mathcal L\mathbf b,
 &
 \mathbf u+c_\sigma\mathbf w&=\beta_\rho\ddot{\mathbf b}.
 \label{eq:op5-sum-identities}
\end{align}
Die Dichte ist minus die Divergenz der ersten Kombination, der Impuls ihre Zeitableitung. Die Spannungskomponenten des Quellentensors sind minus \(\sigma\) angewandt auf die zweite Kombination. Die Kontinuitätsgleichung kürzt so die Dichteableitung gegen die Impulsdivergenz. In der Impulsbilanz kürzt sich \(\beta_\rho\mathcal L\ddot{\mathbf b}\) gegen die Divergenz der Summenspannung. Für \(\mathbf b=\nabla S\) und \(K_L=\lambda+2\mu>0\), mit den Lamé-Koeffizienten \(\lambda,\mu\) der ersten Komponente, wird \(\mathcal L\mathbf b=K_L\nabla\Delta S\). Setzt man wie oben die Ruhedichte des Materials mit seiner Yukawa-Massendichte gleich, wird mit \eqref{eq:op5-density}
\begin{equation}
 \delta\mathcal J_h=-\frac{\delta T^{00}}{v}
                  =\frac{\beta_\rho K_L}{v}\Delta^2 S .
 \label{eq:op5-elas-source}
\end{equation}
Die zwei Laplace-Operatoren in \eqref{eq:op5-elas-source} kommen von der longitudinalen Elastodynamik und der Divergenz in der Dichte. Die Konstruktion mit einem direkten Higgsantrieb setzt \(\beta_\rho\ne0\) voraus. Bei \(\beta_\rho=0\) verschwinden in diesem Ansatz sowohl der mechanische Tensor als auch diese Quelle.

\subsubsection{Wo die Kopplung an das OP5-Feld eingeht}

Das reduzierte OP5-Modell setzt einen gerichteten massiven Vermittler mit Massenparameter \(m>0\) und ein skalares Vergleichsfeld \(h\) an. Für ein glattes Präparationsprofil \(H\) und verschwindende frühere Störungen gilt \(h=d_m H/m^2\) in Fourierdarstellung. Wir verwenden den Ansatz mit unverändertem freiem Einstein-Hauptteil. Für \(x=\omega^2\), \(q=|\mathbf k|^2\), \(a_\perp=|\mathbf k_\perp|^2\) und \(d=q-x\) hat der Antriebsoperator \(\mathcal P_{\rm dyn}\) das Symbol
\begin{align}
 d_m&=m^2+r_mq-x, &
 N_m&=\tfrac23(x^2-a_\perp x+a_\perp^2),\\
 \mathcal P_{\rm dyn}(\omega,\mathbf k)
 &=\frac{[Z(\Omega_0^2+c_h^2q-x)+\tfrac32g_0^2d]d_m-g_m^2N_m}{m^2}.
 \label{eq:op5-drive}
\end{align}
Die Koeffizienten und die gerichtete Quellenform sind vorgegeben. Eine in den Notizen geprüfte Wahl ist \(Z=1\), \(r_m=c_h^2=1/2\), \(g_0^2=g_m^2=3/100\), \(m^2,\Omega_0^2>0\). Sie dient hier als Rechenbeispiel. Liest man das Vergleichsfeld als Radialmodus des Higgsradius von W2, erzwingt W2 \(c_h=1\), bei minimaler Kopplung \(g_0=0\), und mangels gerichteten Vermittlers \(g_m=0\); \(c_h^2=1/2\) ist also eine Vergleichswahl, die W2 nicht trägt (Runde v4.10.0-r33 \cite{layers}). Die Kegelgrenze der Kausalität verlangt, dass beide Vermittler und der Higgsradius langsamer als das Licht laufen. Mit lichtschnellem Higgsradius ist die gerichtete Kopplung für keine Wahl der Vermittlergeschwindigkeiten verträglich (Runden v4.10.0-r47 und v4.10.0-r49 \cite{layers}). Im Vakuum begrenzen Laborgrenzen der Lorentz-Verletzung, die für Fermionen gemessen und auf den Higgs-Sektor übertragen sind, die Kopplung auf \(g_m^2\lesssim10^{-26}\). Der Beispielwert \(g_m^2=3/100\) ist dort ausgeschlossen und nur in einem Material denkbar, dessen Ruhesystem die Präparation ist (Runden v4.10.0-r50 und v4.10.0-r51 \cite{layers}; Anhang~\ref{sec:anhang-b-annahmen}). Auf ein Profil \(\delta H\) antwortet das reduzierte Modell bei passender Präparation auf den Antrieb \(\delta J_{\rm OP5}=\mathcal P_{\rm dyn}\delta H\) \cite{op5-channel}. Die Rechenschritte von hier bis zur Bilanz zweiter Ordnung sind in Anhang~\ref{sec:anhang-b} einzeln hergeleitet und mit Wolfram geprüft.

Die Wahl der zwei Quellenformen hat im reduzierten Modell eine konkrete Aufgabe. Mit \(p^\mu=(\omega,\mathbf k)\) ist \(d=p_\mu p^\mu\). Am Einstein-Kanal wird die identisch transversale Verbesserung \(S^{\mu\nu}=p^\mu p^\nu-\eta^{\mu\nu}d\) eingesetzt. In vier Dimensionen erfüllen ihre Spur und ihre doppelte Kontraktion
\[
 p_\mu S^{\mu\nu}=0,\qquad
 \operatorname{tr}S=-3d,\qquad S:S=3d^2 .
\]
Der spurgespiegelte Austauschterm der beiden Quellen ist proportional zu
\[
 \frac{S:S-(\operatorname{tr}S)^2/2}{d}=-\frac32 d .
\]
Nach Kürzung des Pols bleibt der lokale Beitrag, der in den Antriebsoperator \eqref{eq:op5-drive} eingeht. Auf derselben Verbesserung verschwindet auch die erhaltene Fierz--Pauli-Kontraktion \(S:S-(\operatorname{tr}S)^2/3\). Der gerichtete massive Kontrast in dieser Rechnung benutzt deshalb die eigens gewählte, zur Messachse transversale Quellenform mit dem Zähler \(N_m\). Die gemeinsame Herkunft beider Quellenformen aus einer physischen Wirkung bleibt zu zeigen \cite{op5-channel}. Ein erster Feldansatz gibt den Vermittlerterm mit eindeutigen Kopplungen exakt wieder, setzt aber einen zusätzlichen Skalar (Runde v4.10.0-r38 \cite{layers}). Für sich genommen ist der gerichtete Term ein gewöhnlicher Fierz--Pauli-Austausch einer gesetzten, gerichteten und erhaltenen Quelle; offen ist, ob Rahmen und Achse dieser Quelle zu einem Feld oder zur Auslese gehören (Runde v4.10.0-r45 \cite{layers}).

Die Materierechnung liefert \(\delta\mathcal J_h\). Der Vergleich verlangt \(\delta J_{\rm OP5}\). Wir können beide in deklarierter Quellen-Normierung gleichsetzen und ausrechnen, was daraus folgt. Eine Wirkung, die diese Gleichsetzung aus dem gemeinsamen Träger herleitet, liegt noch nicht vor. Am Higgs-Sektor von W2 folgt sie für den skalaren Kanal, wenn das Vergleichsfeld der Higgs-Radialmodus ist; dann fehlt aber der gerichtete Vermittler, an dem der Kontrast hängt (Runde v4.10.0-r33 \cite{layers}). Gleichung \eqref{eq:op5-match} markiert die zusätzliche Modellannahme; was sie trägt und was ohne sie bestehen bleibt, steht in Anhang~\ref{sec:anhang-b-annahmen}:
\begin{equation}
 \delta J_{\rm OP5}=\delta\mathcal J_h .
 \label{eq:op5-match}
\end{equation}
Mit \(D_0=\partial_t^2-\Delta\) wähle eine glatte, räumlich kompakte Funktion \(V\), die vor der Schaltung null ist, und setze
\begin{equation}
 \delta H=-D_0\Delta^2 V,\qquad
 S=-\frac{v}{\beta_\rho K_L}D_0\mathcal P_{\rm dyn}V .
 \label{eq:op5-profile}
\end{equation}
Die konstantkoeffizientigen Operatoren kommutieren. Beim Einsetzen in \eqref{eq:op5-elas-source} erhält man \(\delta\mathcal J_h=\mathcal P_{\rm dyn}\delta H\). Dieselbe Wahl hat eine zweite Konsequenz. Mit \(\mathbf b_*=(v/(\beta_\rho K_L))\nabla\mathcal P_{\rm dyn}V\) gilt \(\mathbf b=-D_0\mathbf b_*\). Folglich ist \(\delta T=-D_0 T_*\), wobei \(T_*\) der separat erhaltene Tensor aus \eqref{eq:op5-tensor} für \(\mathbf b_*\) ist.

Für die gewöhnliche linearisierte Gravitation verwenden wir harmonische Eichung. Ohne frühere Störung oder einlaufende freie Welle gehört zur Feldgleichung die folgende retardierte Lösung:
\begin{equation}
 D_0\bar\gamma_{\mu\nu}=16\pi G\,\delta T_{\mu\nu},
 \qquad \bar\gamma_{\mu\nu}=-16\pi G\,T_{*,\mu\nu} .
 \label{eq:op5-compact}
\end{equation}
Der gewöhnliche lineare Krümmungswechsel dieses präparierten mechanischen Beitrags verschwindet außerhalb seines räumlichen Trägers, auch während der Schaltung. Die Aussage benutzt den gemeinsamen Wellenfaktor der Quelle und die gewählten Anfangsdaten \cite{op5-motion}.

\subsubsection{Was auf derselben Messachse erhalten bleibt}

Ist damit auch der gerichtete Kontrast verschwunden? Auf dem statischen Plateau wird \(\delta H=\Delta^3 V\). Sei \(\xi=r_\perp^2/R^2\) die radiale Koordinate, \(\chi(\xi)\ge0\) ein glatter kompakter Zuschnitt, der nahe der Achse eins ist, und \(p(z)\) ein glattes Profil mit positivem Integral und Träger zwischen den Anschlüssen. Setze für \(\mathcal A\ne0\)
\begin{equation}
 V=-\frac{\mathcal A R^6}{230400}\,p(z)\chi(\xi)
       (\xi^5-\kappa\xi^6),\qquad
 \kappa=\frac{\int_0^\infty \xi^5\chi(\xi)\,d\xi}
              {\int_0^\infty \xi^6\chi(\xi)\,d\xi}.
 \label{eq:op5-axis-profile}
\end{equation}
Das Querintegral von \(V\) verschwindet je Längsschnitt.

Im flachen Kern kann man den Achsenterm Schritt für Schritt verfolgen. Mit der dimensionslosen Radialkoordinate \(\xi\) gilt \(\Delta_\perp \xi^j=4j^2\xi^{j-1}/R^2\). Dreimaliges Anwenden auf das Polynom aus \eqref{eq:op5-axis-profile} liefert \(230400(\xi^2-4\kappa\xi^3)/R^6\). Enthält ein Term von \(\Delta^3\) eine Längsableitung, bleiben ihm höchstens zwei radiale Laplace-Operatoren. Sein niedrigster radialer Grad ist daher \(\xi^3\). In der Nähe der Achse gelten für \(p(z)\ne0\) die normierten Formen
\begin{equation}
 \frac{\delta H}{-\mathcal A p(z)}=\xi^2+O(\xi^3),\qquad
 \frac{\Delta_\perp\delta H}{-16\mathcal A p(z)/R^2}=\xi+O(\xi^2).
 \label{eq:op5-normalized-axis}
\end{equation}
Die Profiländerung selbst ist auf der Achse null. Der massive Vergleich benutzt neben dem skalaren Feld \(h=d_mH/m^2\) das Feld \(f=\Delta_\perp H/m^2\), das aus demselben Profil mit dem transversalen Laplace-Operator gebildet wird. Nach \eqref{eq:op5-normalized-axis} ist auch sein Eintrag \(\delta f(0,z)\) auf der Achse null. Die zweite transversale Ableitung dieses Eintrags kann trotzdem von null verschieden sein \cite{op5-motion}.

\begin{figure}[!htbp]
\centering
\begin{tikzpicture}[>=Latex,x=43cm,y=13.4cm,font=\small]
  \draw[qs axis,->] (0,0)--(.215,0) node[right] {\(\xi=r_\perp^2/R^2\)};
  \draw[qs axis,->] (0,0)--(0,.211);
  \foreach \x in {.05,.1,.15,.2} {\draw[qs axis] (\x,-.004)--(\x,.004);}
  \node[anchor=north east] at (0,0) {\(0\)};
  \node[below=3pt] at (.1,0) {\(0{,}1\)};
  \node[below=3pt] at (.2,0) {\(0{,}2\)};
  \foreach \y in {.1,.2} {\draw[qs axis] (-.0012,\y)--(.0012,\y);}
  \node[left=4pt] at (0,.1) {\(0{,}1\)};
  \node[left=4pt] at (0,.2) {\(0{,}2\)};
  \draw[qs rust,very thick,dashed,domain=0:.19,samples=2] plot (\x,{\x});
  \draw[qs blue,very thick,domain=0:.19,samples=38] plot (\x,{\x*\x});
  \node[anchor=south east,qs rust] at (.177,.179) {\(\xi\)};
  \node[anchor=south west,qs blue] at (.163,.038) {\(\xi^2\)};
\end{tikzpicture}
\caption{Führende Terme nahe der Achse. Für \(p(z)\ne0\) beginnt das normierte Profil \(\delta H/(-\mathcal A p)\) mit \(\xi^2\), und sein normierter transversaler Laplace-Ausdruck \(\Delta_\perp\delta H/(-16\mathcal A p/R^2)\) mit \(\xi\). Damit verschwinden auf der Achse beide Werte, doch die zweite transversale Ableitung der Laplace-Antwort liefert den Gezeitenkontrast. Gezeigt werden die führenden Terme.}
\label{fig:op5-profil}
\end{figure}


In der geometrischen Optik ergibt sich für eine transversale Einheitsrichtung \(e\) auf derselben Achse der Gezeitenkontrast
\begin{equation}
 \Delta\mathcal I=-\frac12\int\partial_e^2(\delta f)(0,z)\,dz
                 =\frac{16\mathcal A}{m^2R^4}\int p(z)\,dz .
 \label{eq:op5-contrast}
\end{equation}
Der Ausdruck ist ungleich null, obwohl der in \eqref{eq:op5-compact} betrachtete gewöhnliche lineare Außenbeitrag kompakt ist. In der idealisierten Krümmungsauslese multipliziert eine gesetzte Differenz der effektiven Kopplungen der beiden Auslesekanäle diesen Kontrast. Erst mit einer solchen Kopplungsdifferenz und kalibrierten Anschlüssen lässt sich daraus eine Laufzeitdifferenz für einen Messaufbau bestimmen \cite{op5-axis,op5-motion}. Von der Ausleseseite gelesen ist die gerichtete Kopplung die kovariante Auslese des Gezeitenfelds quer zur Messachse eines Detektors mit Ruhesystem \(u\) und Achse \(n\) (Runde v4.10.0-r46 \cite{layers}). Rahmen und Achse wählt diese Auslese nicht selbst aus; sie sind eine Setzung des Detektors (Runde v4.10.0-r54 \cite{layers}).

Der Kontrast \(\Delta\mathcal I\) ist ein Integral einer Krümmungskomponente des Vergleichsfelds für eine feste Polarisation. Ein statischer konformer Term geht auf der Achse als integrierte Randableitung ein, und ein rein longitudinaler Anteil hat keine lineare Krümmung. Die Quelle kann so an den Enden identische lokale Vergleichsfelder liefern und auf dem Weg trotzdem den gerichteten Kontrast ändern \cite{op5-axis}. Für die dynamisch präparierte Profilfamilie liefert die Azimutalmittelung über ein rotationssymmetrisch gewichtetes, achsenparalleles Bündel mit konstantem Querschnitt im flachen Kern den Faktor
\begin{equation}
 1-36\kappa\langle\xi\rangle
 \label{eq:op5-beam}
\end{equation}
zum auf der Achse berechneten reinen Riemann-Kontrast. In dieser Mittelung verschwinden die Längsableitungsbeiträge nach der Integration in \(z\). Der Faktor ist eine Kontrolle der Näherung am statischen Plateau. Bei Beugung, wechselndem Querschnitt oder Durchlaufen der Strecke während des Umschaltens ist die optische Übertragung eine weitere Rechnung \cite{op5-motion,op5-axis}.

\paragraph{Planare Kontrolle.}
Beim planaren statischen Vergleich mit dem Faltungskern \(K_m(z)=m e^{-m|z|}/2\) und einer im Inneren der Strecke getragenen Dichtedifferenz \(\delta\rho\) gilt für \(f=K_m*\delta\rho\)
\begin{equation}
 \int_0^L f(z)\,dz
   =\int_0^L\delta\rho(z)\,dz-\frac{f(0)+f(L)}{m}.
 \label{eq:op5-planar-check}
\end{equation}
Die Faltung erhält das Integral über die ganze Gerade. Außerhalb des Trägers sind ihre beiden Ausläufer exponentiell und ihre Integrale werden durch die beiden Endwerte bestimmt. Bei gleicher integrierter Dichte und gleichen Endwerten in den beiden Stellungen bleibt damit aus diesem planaren Beitrag keine führende Differenz. In der axialen Konstruktion von \eqref{eq:op5-contrast} misst dieselbe geometrische Strecke eine transversale Krümmungsableitung. Gleiches Querintegral und gleiche lokale Endwerte schließen diesen Kontrast nicht aus. Das ist die Kontrolle, die den Wechsel des Profils mit festem Messweg motiviert \cite{op5-axis}.

\subsubsection{Warum die zweite Ordnung eine eigene Bilanz braucht}

\paragraph{Arbeit beim Schalten.}
Die innere Feder überträgt Impuls zwischen den beiden Materialien. Ihre vorgeschriebene Ruhelage \(\mathbf d\) ist zugleich eine Arbeitsvariable. Setze \(\mathbf R_{\rm F}=\mathbf u-\mathbf w-\mathbf d\), \(\sigma_1=\sigma(\mathbf u)\) und \(\sigma_2=c_\sigma\,\sigma(\mathbf w)\). Für die zur zweiten Ordnung beitragende mechanische Energiedichte gilt
\begin{align}
 \partial_t e_{\rm mech}&=\nabla\cdot\mathbf F+P_{\rm Schalt},\\
 \mathbf F&=\sigma_1\dot{\mathbf u}+\sigma_2\dot{\mathbf w},
 &P_{\rm Schalt}&=-k\,\mathbf R_{\rm F}\cdot\dot{\mathbf d}.
 \label{eq:op5-drive-energy}
\end{align}
Die Vorzeichen folgen aus der Federenergie \(k|\mathbf R_{\rm F}|^2/2\). Beim Differenzieren hebt sich die Arbeit der beiden inneren Federkräfte gegen den Anteil mit \(\dot{\mathbf u}-\dot{\mathbf w}\) auf. Übrig bleibt der von außen vorgegebene Anteil mit \(\dot{\mathbf d}\). Ein ortsfester Speicher mit \(\partial_t e_{\rm res}=-P_{\rm Schalt}\) gleicht die Arbeit im Volumen aus. Aus der Summe folgt auf dem Plateau noch der Dichtebeitrag \(\nabla\cdot\int\mathbf F\,dt\). Sein Volumenintegral verschwindet für kompakten Fluss. Höhere Momente können erhalten bleiben \cite{op5-feder}.

\paragraph{Welche Verschiebung am Plateau fehlt.}

Die lineare Gegenbewegung ist ein Befund für die präparierte Quelle. Um die zweite Ordnung zu prüfen, betrachten die Forschungsnotizen dieselben beiden Hauptkomponenten auf dem Plateau mit einem Saint-Venant--Kirchhoff-Gesetz und ausschließlich innerer Paarfeder. Sei \(\mathsf H=\nabla_X\mathbf u_1\) die symmetrische erste Deformation im Referenzchart, mit \(\mathbf u_1=\nabla_X\varphi\) und \(\mathbf w_1=-\mathbf u_1/c_\sigma\). Für die zweite Summenverschiebung \(S_2=\mathbf u_2+c_\sigma\mathbf w_2\) hat die Summe der Piola-Tensoren den Koeffizienten
\begin{align}
 (P_1+P_2)^{(2)}&=\mathsf C\,\varepsilon(S_2)
                     +(1+1/c_\sigma)N(\mathsf H),\\
 N(\mathsf H)&=\tfrac{\lambda}{2}\operatorname{tr}(\mathsf H^2)I
                 +3\mu\mathsf H^2+\lambda(\operatorname{tr}\mathsf H)\mathsf H .
\end{align}
Für \(\mu>0\), \(3\lambda+2\mu>0\) und eine nichttriviale kompakte longitudinale erste Verschiebung folgt
\begin{equation}
 \int\operatorname{tr}N(\mathsf H)\,d^3X
 =\left(\tfrac52\lambda+3\mu\right)\int|\mathsf H|^2\,d^3X>0 .
 \label{eq:op5-second}
\end{equation}

Bei kompaktem \(S_2\) erzwingt die statische Summengleichung nach Testen mit \(X\) eine verschwindende integrierte Spur. Der lineare Dehnungsterm integriert sich dabei zu null. Der positive quadratische Beitrag widerspricht dieser notwendigen Bedingung. Damit scheitert eine kompakte zweite Summenverschiebung unter den angegebenen Material- und Halterungsannahmen \cite{op5-elastic-limit}.

Die Spuridentität lässt sich im selben Referenzchart prüfen. Für die Hesse-Matrix \(\mathsf H\) von \(\varphi\) mit kompakter erster Verschiebung ist
\[
 (\Delta\varphi)^2-\operatorname{tr}(\mathsf H^2)
 =\nabla_X\cdot\big[(\Delta\varphi)\nabla_X\varphi
                    -\mathsf H\nabla_X\varphi\big].
\]
Der Randterm verschwindet. In der integrierten Spur von \(N\) lassen sich daher \((\operatorname{tr}\mathsf H)^2\) und \(\operatorname{tr}\mathsf H^2\) zusammenfassen. Damit erhält man den Koeffizienten aus \eqref{eq:op5-second}. Das Argument prüft eine mechanische Summenverschiebung und verwendet keine optische Auslese. Es macht auch deutlich, welche Größen getrennt bilanziert werden müssen. Der zusätzliche Term im Restmassentransport ist
\begin{equation}
 -\nabla\cdot(\rho_1\mathbf u_2+\rho_2\mathbf w_2).
 \label{eq:op5-second-density}
\end{equation}
Die statische Summengleichung in \(X\) bestimmt die Kombination mit \(c_\sigma\). Der gravitative Beitrag aus \eqref{eq:op5-second-density} braucht die mit den Ruhedichten gewichtete Kombination. Für den räumlichen Stress ist außerdem der Cauchy-Tensor am selben Raumpunkt zu bilden. Die Federkräfte setzen dort jeweils an unterschiedlichen Endpunkten an und benötigen eine Bilanz des Verbindungsstresses. Bei nichtzentralen Kräften kommt die Drehmomentbilanz der Halterung hinzu \cite{op5-elastic-limit}.

Der vorhandene Reservoirabgleich kompensiert die bereits berechneten direkten Dichteanteile auf dem Plateau. Der vollständige gewöhnliche Beitrag zweiter Ordnung bleibt zu bestimmen \cite{op5-reservoir,op5-elastic-limit}.

\subsubsection{Wie sich eine Laufzeit von der Anschlussphase trennen ließe}
\label{sec:op5-calibration}

Die Krümmung in \eqref{eq:op5-contrast} ist eine Größe im Modell. Wie sähe eine Auslese aus, die sie von einer bloßen Änderung der Anschlüsse trennen könnte? Betrachte dazu eine kohärente, verlustfreie Zweimodenstrecke in der linearen Einweg-Näherung. Nach dem Ausklammern der gemeinsamen Grundlaufzeit seien \(M_A=LF\) und \(M_B=LUF\). \(L\) und \(F\) sind invertierbare Empfangs- und Sendematrizen. \(U\) ist der unitäre Streckenkontrast zwischen den Profilstellungen. Wenn die beiden Enden in beiden Stellungen dieselben Matrizen haben, gilt
\begin{equation}
 M_A^{-1}M_B=F^{-1}UF .
 \label{eq:op5-ab-calibration}
\end{equation}
Die Eigenphasen sind dann diejenigen der Strecke. Diese Aussage benutzt die Ähnlichkeit und braucht keine diagonalen Anschlüsse. Falls sich die Enden mit dem Profil ändern, kann ein profilsynchroner Phasenterm am Sender denselben Frequenzgang wie ein mit der Laufzeitmatrix kommutierender Zusatz auf der Strecke erzeugen. Der Quotient allein entscheidet in diesem Fall nicht über den Ort des Effekts \cite{op5-budget}.

Ein unabhängiger Nullweg mit bekanntem Einheitskontrast würde im Idealfall die gleichen aktuellen Endmatrizen teilen. Für relative Abweichungen \(S,R\) am Empfangs- und Sendeende lässt sich diese Forderung direkt schreiben. Ist \(C=LF\) die Nullmessung und \(M=LSURF\) die Messung, so ist \(C^{-1}M\) für invertierbares \(S\) ähnlich zu \(U(RS)\). Reziproke Abweichungen mit \(RS=I\) erhalten das Streckenspektrum. Bei \(RS=c_0 I\), \(c_0\ne0\), erhalten beide Eigenwerte denselben Faktor, sodass ihre korrekt entfaltete Phasendifferenz gleich bleibt. Alternativ könnte eine im Messzweig unabhängig begründete Nullstellung \(M_0=LSRF\) die gleichen relativen Anschlüsse wie die Messung verwenden. Dann gilt \(M_0^{-1}M=F^{-1}R^{-1}URF\). Die Voraussetzung ist in beiden Fällen eine kohärent geteilte Endpunktantwort. Die axiale Nullstellung des Feldvergleichs liefert dafür ein Kriterium an das Profil. Sie bestimmt noch nicht den Abgleich eines physischen Messanschlusses \cite{op5-nullpfad,op5-axis}.

Auch bei abgeglichenen Enden bleibt eine Anforderung an die Stabilität der Phase. Schreibe die fehlerbehafteten Matrizen in der Feldbasis als \(\widehat M_A=L(I+N_A)F\) und \(\widehat M_B=L(U+N_B+D_{\rm E})F\), mit \(D_{\rm E}\) als Änderung der Endpunktantwort in Stellung B. Für invertierbares \(I+N_A\) ist der Messquotient \(\widehat M_A^{-1}\widehat M_B\) ähnlich zu \(Q=(I+N_A)^{-1}(U+N_B+D_{\rm E})\). Ihr Fehler gegenüber \(U\) ist exakt
\begin{equation}
 Q-U=(I+N_A)^{-1}(N_B-N_AU+D_{\rm E}) .
 \label{eq:op5-phase-error}
\end{equation}
Bezeichne \(r\) eine Schranke für \(\|Q-U\|_2\). Für normales \(U\) liegt jeder gestörte Eigenwert höchstens im Abstand \(r\) zu einem idealen Eigenwert. Sind die Kreise disjunkt, bleibt die Zuordnung erhalten. Für \(r<1\) beträgt die Winkelabweichung jedes zugeordneten Eigenwerts höchstens \(\arcsin r\). Über zwei Frequenzen mit Abstand \(\Delta\omega\) erhält man für die Sekante der differentiellen Phase nach korrekter Entfaltung die Fehlergrenze \(4\arcsin r/\Delta\omega\). Die vier Beiträge stammen von zwei Eigenphasen an jeweils zwei Frequenzen \cite{op5-budget}. Die Kalibrieraufgabe erhält damit eine präzise Form. Das am Modell vorhergesagte Signal muss sich sowohl gegen die verbleibende Endpunktänderung als auch gegen die Phasenfehler behaupten.

Die Fallstudie zeigt damit, welche Schritte sich gemeinsam rechnen lassen. Die W2-Massenparameter bestimmen die radiale Materiequelle. Für das präparierte Profil sind gewöhnliches lineares Außenfeld und gerichteter Achsenkontrast getrennt berechenbar. Der gewöhnliche Teil folgt allein aus der linearisierten Einstein-Gleichung. Der gerichtete Teil hängt an der Kopplungsannahme \eqref{eq:op5-match} und an einem gerichteten, unterluminalen massiven Vermittler, den W2 nicht enthält. Mit dem lichtschnellen Higgsradius von W2 verletzt die gerichtete Kopplung die Kegelgrenze, und im Vakuum bleibt nur \(g_m^2\lesssim10^{-26}\); ein beobachtbarer Kontrast verlangt ein Material. Die native Massenauslese, eine Wirkung für Vermittler und Kopplung und die kalibrierte Auslese bleiben weitere Schritte \cite{op5}; OP5 bleibt offen. Den Stand der Vermittlersuche fasst Anhang~\ref{sec:anhang-b-annahmen} zusammen, die Einordnung unter die übrigen offenen Probleme Runde v4.10.0-r55 \cite{crossmap}.
